\section{Introduction}

Choosing an appropriate inductive bias is crucial for generalization \cite{baxter2000model}.
By constraining the hypothesis space to reflect task structure, it can enable effective learning from limited data.
Examples include local connectivity and spatial weight sharing in convolutional networks \cite{lecun1998gradient}, and recurrent state updates with parameters shared across time \cite{elman1990finding}.
Suitable parameter-sharing patterns also enforce equivariance to task-relevant transformations \cite{ravanbakhsh2017equivariance}.

Selecting suitable structural biases by hand requires domain expertise, particularly when the relevant properties of a new task are not known in advance \cite{yeh2022equivariance}.
This motivates learning such biases from data.
For example, \citet{zhou2021meta} meta-learn layer reparameterizations to discover shared parameter structures, including symmetry-inducing patterns.
\citet{atanov2019deep} learn generative priors from convolutional filters of previously trained networks and use them to guide learning on new datasets.
\citet{sun2020sparse} extract task-specific subnetworks by iterative pruning and share parameters where these subnetworks overlap.
\citet{ando2005framework} learn a shared predictive structure across tasks while allowing task-specific predictors.

More generally, a single task may admit multiple near-optimal structures, and its data alone need not identify which constraints will transfer to other tasks.
\citet{baxter2000model} provides a theoretical foundation for learning such constraints across an environment of related tasks.
Under capacity and sampling assumptions, his bounds relate empirical performance across sufficiently many source tasks to the expected risk of the best task-specific hypotheses in a selected space on new tasks from the same environment, with high probability.
Our focus is on extracting one transferable connectivity pattern from diverse trained solutions: a mask that remains useful when fresh weights are learned for a new task.

In this paper, we present a generative approach to learning a shared connectivity mask for related tasks.
Building on generative modeling of trained networks \cite{atanov2019deep,schurholt2022hyper}, we construct a bank of functional maps for each task from networks trained at different sparsity levels.
Each map records neuron responses and connection sensitivities on fixed probe examples, together with the connectivity mask.
Task-specific generators use these banks to propose new masks, extending the search beyond the stored solutions.
A shared, task-conditioned evaluator predicts their loss and guides the search, following the principle of surrogate-guided architecture search \cite{liu2018progressive}.

The order of stored solutions and hidden neurons is arbitrary, so the generators process them as sets \cite{zaheer2017deep}.
Their Transformer encoders omit positional encodings for these elements and average neuron representations to summarize each solution.
An additional penalty encourages identical outputs when the bank is reordered while keeping the sampled noise fixed.
An agreement loss encourages generators for different tasks to propose similar masks after aligning hidden neurons.

Search alternates mask generation with candidate evaluation on all source tasks: fresh weights are trained on support data, and loss is measured on separate query data.
These measurements update the evaluator and generators, while selected trained models replenish the functional banks.
Finally, a criterion based on the worst source-task performance selects one mask with a fixed connection budget using separate selection data.
We freeze this mask before evaluating unseen tasks, learning weights independently for each task.

Our main contributions are as follows:
\begin{itemize}
    \item We propose a generative framework for learning a shared connectivity mask as a structural inductive bias, with a fixed connection budget and independently trained task-specific weights.
    \item We introduce functional maps of trained networks as inputs to mask generation, combining neuron responses, connection sensitivities, and connectivity with an encoder invariant to solution and hidden-neuron ordering.
    \item We develop a cooperative search procedure that combines task-specific generators, agreement between their proposals, and a shared loss evaluator updated through candidate training on source tasks.
    % TODO: Intended empirical contribution, pending final experiments.
    % Confirm fixed-mask transfer across sample sizes and comparisons with
    % dense and matched-budget sparse baselines before retaining this claim.
    \item We demonstrate improved data efficiency on previously unseen tasks using synthetic pattern-detection problems with known reference connectivity and image-based set-regression tasks, with comparisons against dense networks and matched-budget sparse baselines.
\end{itemize}
